\documentclass[11pt,a4paper]{article}
\usepackage[T2A]{fontenc}
\usepackage[utf8]{inputenc}
\usepackage[russian]{babel}
\usepackage{amsmath,amssymb,amsfonts}
\usepackage{geometry}
\geometry{top=20mm,bottom=20mm,left=20mm,right=20mm}
\usepackage{graphicx}
\usepackage{booktabs}
\usepackage{tabularx}
\usepackage{array}
\usepackage{xcolor}
\usepackage{hyperref}
\usepackage{listings}

% Цветовая палитра оформления
\definecolor{primary}{RGB}{20, 50, 110}
\definecolor{secondary}{RGB}{0, 102, 102}
\definecolor{accent}{RGB}{160, 40, 40}
\definecolor{codebg}{RGB}{248, 249, 250}
\definecolor{boxbg}{RGB}{240, 244, 250}
\definecolor{boxframe}{RGB}{70, 100, 160}
\definecolor{warnbg}{RGB}{255, 248, 240}
\definecolor{warnframe}{RGB}{210, 120, 30}
\definecolor{darkgreen}{RGB}{20, 120, 40}

\hypersetup{
    colorlinks=true,
    linkcolor=primary,
    citecolor=secondary,
    urlcolor=accent,
    pdftitle={Реализация и вычислительные приемы метода Монте-Карло},
    pdfauthor={МГТУ им. Н.Э. Баумана, каф. ФН-4}
}

% Настройка листингов кода C++
\lstset{
    language=C++,
    backgroundcolor=\color{codebg},
    basicstyle=\ttfamily\footnotesize,
    keywordstyle=\color{primary}\bfseries,
    commentstyle=\color{darkgreen}\itshape,
    stringstyle=\color{accent},
    numbers=left,
    numberstyle=\tiny\color{gray},
    stepnumber=1,
    numbersep=8pt,
    tabsize=4,
    showspaces=false,
    showstringspaces=false,
    breaklines=true,
    frame=single,
    rulecolor=\color{boxframe},
    captionpos=b
}

\newcommand{\infobox}[2]{%
  \par\vspace{3mm}\noindent
  \fcolorbox{boxframe}{boxbg}{%
    \parbox{\dimexpr\textwidth-2\fboxsep-2\fboxrule\relax}{%
      \textbf{\color{primary}#1}\vspace{1.5mm}\par\small #2
    }%
  }\par\vspace{3mm}
}

\newcommand{\warnbox}[2]{%
  \par\vspace{3mm}\noindent
  \fcolorbox{warnframe}{warnbg}{%
    \parbox{\dimexpr\textwidth-2\fboxsep-2\fboxrule\relax}{%
      \textbf{\color{accent}#1}\vspace{1.5mm}\par\small #2
    }%
  }\par\vspace{3mm}
}

\title{\vspace{-10mm}\Huge\bfseries\color{primary} Архитектура реализации и вычислительные приемы в лабораторной работе №6\\[3mm]
\Large\color{secondary} Разбор программного комплекса на C++20, численные алгоритмы, математические хитрости и оптимизации}
\author{\textbf{Кафедра «Физика» (ФН-4), МГТУ им. Н.Э. Баумана}\\
Курс «Вычислительная физика» (7 семестр)}
\date{2026}

\begin{document}

\maketitle

\begin{abstract}
\noindent В настоящем документе детально описана программная и алгоритмическая реализация методов интегрирования Монте-Карло, разработанных в рамках лабораторной работы №6. Разобраны структура заголовочного модуля \texttt{monte\_carlo.hpp} и управляющей программы \texttt{lab6.cpp}. Освещены ключевые проблемы вычислительной математики, возникающие при моделировании больших выборок ($N = 10^6 \dots 10^7$ точек), и показано, какими алгоритмическими приемами («хитростями») они были решены.

В частности, подробно рассмотрены: преодоление катастрофической потери точности машинной арифметики при помощи однопроходного рекуррентного алгоритма Велфорда (Welford's algorithm); тонкости 64-битного генератора псевдослучайных чисел \texttt{std::mt19937\_64}; честный учет вычислительной сложности антитетических пар; построение оптимальных контрольных функций (почему двухточечная хорда превосходит разложение Тейлора в 10 раз, а параболическая контрольная функция снижает дисперсию в 250 раз); защита от возникновения нечисловых значений (\texttt{NaN}) при вычислении интеграла с вертикальной производной $\sqrt{3 - x^2}$; сопоставление со сбоями формулы Симпсона из лабораторной работы №5 (Вариант 7).
\end{abstract}

\tableofcontents
\newpage

\section{Общая архитектура программного комплекса}

Программная часть лабораторной работы №6 спроектирована по модульному принципу в соответствии с современными стандартами языка C++20:
\begin{enumerate}
    \item \textbf{Библиотечный заголовочный модуль \texttt{monte\_carlo.hpp}:}  
    Содержит в себе реализацию математических примитивов, структур данных, алгоритмов генерации чисел, пяти методов интегрирования Монте-Карло и классических формул квадратур Ньютона--Котеса (для сквозного сопоставления с результатами Лабораторной работы №5).
    
    \item \textbf{Управляющий расчетный файл \texttt{lab6.cpp}:}  
    Выполняет вычислительный эксперимент:
    \begin{itemize}
        \item Расчет базового интеграла задания: $I_1 = \int_0^1 e^x dx = e - 1 \approx 1.718281828459$;
        \item Расчет дополнительного интеграла задания: $I_2 = \int_0^{\pi/2} \sin(x) dx = 1.000000000000$;
        \item Расчет интеграла Варианта 7 из Лаб №5: $I_{circ} = \int_0^{\sqrt{3}} \sqrt{3 - x^2} dx = \frac{3\pi}{4} \approx 2.356194490192$;
        \item Построение таблиц сходимости при $N \in \{10^2, 10^3, 10^4, 10^5, 10^6, 10^7\}$;
        \item Экспорт всей числовой информации в формате CSV (\texttt{lab6\_results.csv}).
    \end{itemize}
    
    \item \textbf{Скрипт визуализации \texttt{plot\_results.gp}:}  
    Автоматизирует генерацию высококачественных графиков в векторном растре PNG с использованием пакета Gnuplot.
\end{enumerate}

\begin{figure}[h!]
\centering
\fbox{
\begin{minipage}{0.92\textwidth}
\textbf{Взаимосвязь модулей проекта:}
\begin{center}
\texttt{monte\_carlo.hpp} \ $\longrightarrow$ \ \texttt{lab6.cpp} \ $\xrightarrow{\text{g++ -O3}}$ \ \texttt{lab6.exe} \\
$\downarrow$ \\
\texttt{lab6\_results.csv} \ $\longrightarrow$ \ \texttt{plot\_results.gp} \ $\xrightarrow{\text{gnuplot}}$ \ \texttt{fig\_*.png}
\end{center}
\end{minipage}
}
\end{figure}

---

\section{Хитрость 1. Преодоление катастрофической потери точности (Алгоритм Велфорда)}

\subsection{В чем состоит скрытая проблема наивного подхода?}
В учебниках часто приводится школьная формула выборочной дисперсии:
\begin{equation}
    s^2 = \frac{1}{N - 1} \sum_{i=1}^N (x_i - \bar{x})^2 = \frac{1}{N - 1} \left[ \sum_{i=1}^N x_i^2 - \frac{1}{N} \left(\sum_{i=1}^N x_i\right)^2 \right].
    \label{eq:naive_var}
\end{equation}

На первый взгляд, формула \eqref{eq:naive_var} удобна тем, что позволяет найти дисперсию за один проход: достаточно параллельно накапливать сумму значений $S_1 = \sum x_i$ и сумму квадратов $S_2 = \sum x_i^2$.

\warnbox{Опасность наивной формулы (Catastrophic Cancellation)}{
В стандарте IEEE 754 число двойной точности (\texttt{double}) имеет 53 бита мантиссы, что обеспечивает около 15--17 значащих десятичных цифр.

Когда число испытаний велико ($N = 10^7 \dots 10^8$), обе величины $S_2$ и $\frac{S_1^2}{N}$ становятся огромными числами порядка $O(N)$. При этом их разность $S_2 - \frac{S_1^2}{N} = (N - 1)s^2$ имеет порядок $O(N \sigma^2)$.

Если дисперсия мала (например, при использовании методов понижения дисперсии $\sigma^2 \sim 10^{-4}$), первые 8--10 цифр у чисел $S_2$ и $\frac{S_1^2}{N}$ \textbf{полностью совпадают}. В результате их вычитания значащие цифры уничтожаются, и мы вычитаем два близких гигантских числа с плавающей точкой. Это явление в вычислительной математике называется \textbf{катастрофической потерей точности}. На практике это приводит к тому, что вычисленная дисперсия теряет до 10 порядков точности, а в некоторых случаях из-за машинного округления разность оказывается отрицательной, что влечет извлечение корня из отрицательного числа и крах программы (\texttt{NaN})!
}

\subsection{Математическое решение: Рекуррентный алгоритм Велфорда (1962)}
Чтобы получить абсолютно численно стабильную оценку за один проход без сохранения в памяти всей выборки, в модуле \texttt{monte\_carlo.hpp} реализован алгоритм Б.П.~Велфорда.

Пусть на шаге $k$ уже известны среднее $\bar{x}_{k-1}$ и величина $M_{2, k-1} = \sum_{i=1}^{k-1} (x_i - \bar{x}_{k-1})^2$. При поступлении нового элемента $x_k$ пересчет осуществляется по формулам:
\begin{align}
    \delta_k &= x_k - \bar{x}_{k-1}, \label{eq:welford_delta} \\
    \bar{x}_k &= \bar{x}_{k-1} + \frac{\delta_k}{k}, \label{eq:welford_mean} \\
    \delta_{2, k} &= x_k - \bar{x}_k, \\
    M_{2, k} &= M_{2, k-1} + \delta_k \cdot \delta_{2, k}. \label{eq:welford_m2}
\end{align}

По завершении $N$ шагов несмещенная дисперсия находится тривиальным делением без вычитания больших сумм:
\begin{equation}
    s^2 = \frac{M_{2, N}}{N - 1}.
\end{equation}

\subsection{Реализация аккумулятора в коде C++}
В \texttt{monte\_carlo.hpp} этот алгоритм оформлен в виде легковесной структуры \texttt{WelfordAccumulator}:

\begin{lstlisting}
struct WelfordAccumulator {
    uint64_t count = 0;
    double mean = 0.0;
    double M2 = 0.0; // Sum of squared differences from current mean

    void update(double x) {
        count++;
        double delta = x - mean;
        mean += delta / count;
        double delta2 = x - mean;
        M2 += delta * delta2;
    }

    [[nodiscard]] double get_mean() const { 
        return mean; 
    }

    [[nodiscard]] double get_variance() const { 
        return (count > 1) ? (M2 / (count - 1)) : 0.0; 
    }

    [[nodiscard]] double get_std_error() const { 
        return (count > 0) ? std::sqrt(get_variance() / count) : 0.0; 
    }

    [[nodiscard]] double get_ci95_halfwidth() const { 
        // 95% two-sided normal quantile: 1.95996398454
        return 1.95996398454 * get_std_error(); 
    }
};
\end{lstlisting}

\infobox{Преимущества реализации}{
1. Требует строго $O(1)$ памяти (всего 24 байта: один \texttt{uint64\_t} и два \texttt{double}).\\
2. Гарантирует абсолютную точность: $M_2 \ge 0$ строго всегда, появление отрицательной дисперсии математически невозможно.\\
3. Позволяет рассчитывать доверительный интервал «на лету» без второго прохода по данным.
}

\newpage

\section{Хитрость 2. Правильный выбор генератора псевдослучайных чисел (ГПСЧ)}

\subsection{Почему категорически нельзя использовать \texttt{rand()}}
В старых программах на C/C++ часто встречается вызов \texttt{rand() / (double)RAND\_MAX}. Для серьезного моделирования методом Монте-Карло это недопустимо:
\begin{itemize}
    \item \textbf{Малый период:} В Windows CRT константа \texttt{RAND\_MAX} равна всего $32767$ ($2^{15}-1$). При $N = 10^6$ чисел генератор успевает пройти свой полный цикл более 30 раз, выдавая циклически повторяющиеся последовательности!
    \item \textbf{Эффект Марсальи (гиперплоскости):} Линейные конгруэнтные генераторы (LCG) страдают концентрацией $d$-мерных точек на параллельных гиперплоскостях. Если моделировать 2D или 3D блуждания, точки ложатся на регулярные решетки, уничтожая стохастическую природу метода.
\end{itemize}

\subsection{Выбор: 64-битный вихрь Мерсенна (\texttt{std::mt19937\_64})}
В модуле \texttt{monte\_carlo.hpp} создан класс-обертка \texttt{RandomGenerator}:

\begin{lstlisting}
class RandomGenerator {
private:
    std::mt19937_64 rng_;
    std::uniform_real_distribution<double> dist01_{0.0, 1.0};

public:
    explicit RandomGenerator(uint64_t seed = 123456789ULL) : rng_(seed) {}

    void set_seed(uint64_t seed) { rng_.seed(seed); }

    void seed_random_device() {
        std::random_device rd;
        // Combine two 32-bit hardware entropy words into a full 64-bit seed
        rng_.seed((static_cast<uint64_t>(rd()) << 32) ^ rd());
    }

    [[nodiscard]] double uniform01() { 
        return dist01_(rng_); 
    }

    [[nodiscard]] double uniform(double a, double b) { 
        return a + (b - a) * dist01_(rng_); 
    }
};
\end{lstlisting}

\infobox{Особенности генератора}{
1. \textbf{Период:} Составляет $2^{19937} - 1 \approx 4.3 \times 10^{6001}$ (исчерпать такой период невозможно физически).\\
2. \textbf{Равномерность:} Гарантируется равномерное распределение векторов вплоть до размерности $d = 623$.\\
3. \textbf{Аппаратная энтропия:} Метод \texttt{seed\_random\_device()} задействует аппаратный генератор случайных чисел процессора (\texttt{RDRAND}), исключая предсказуемость семян. При этом для научной воспроизводимости результатов в бенчмарках задается фиксированный \texttt{seed = 42}.
}

\newpage

\section{Хитрость 3. Честное сопоставление вычислений для антитетических переменных}

Метод противоположных (антитетических) переменных основан на использовании парных точек $U$ и $1 - U$. 

\subsection{Проблема честного нормирования}
Если мы генерируем $N$ пар $(U_i, 1 - U_i)$, то для каждой пары функция вычисляется дважды: $f(a + (b-a)U_i)$ и $f(a + (b-a)(1 - U_i))$. Таким образом, суммарное число вызовов функции составляет $2N$. 

Если не скорректировать объем выборки, сравнение с обычным методом Монте-Карло при фиксированном $N$ будет нечестным, так как антитетический метод потратит вдвое больше ресурсов процессора!

\subsection{Как это решено в нашем коде}
В функции \texttt{mc\_antithetic\_variates} общее число вызовов функции $N_{evals}$ строго зафиксировано:
\begin{enumerate}
    \item Вычисляется число пар: $N_{pairs} = N / 2$.
    \item Оценивается среднее на одну пару:
    \begin{equation}
        \theta_i = \frac{b - a}{2} \left[ f\big(a + (b - a)U_i\big) + f\big(a + (b - a)(1 - U_i)\big) \right].
    \end{equation}
    \item Аккумулятор Велфорда накапливает дисперсию элементарной оценки пары: $s_{pair}^2 = \mathrm{Var}(\theta)$.
    \item Дисперсия среднего по выборке из $N_{pairs}$ равна:
    \begin{equation}
        \mathrm{Var}(\bar{\theta}) = \frac{s_{pair}^2}{N_{pairs}} = \frac{s_{pair}^2}{N / 2} = \frac{2 \cdot s_{pair}^2}{N}.
    \end{equation}
    \item Для честного сопоставления с базовым методом Монте-Карло (где $\mathrm{Var}(\bar{I}) = \frac{s_{single}^2}{N}$) в результирующей структуре выводится \textbf{эффективная дисперсия на одно вычисление функции}:
    \begin{equation}
        s_{eff}^2 = 2 \cdot s_{pair}^2.
    \end{equation}
\end{enumerate}

Благодаря этому подходу в наших отчетах столбец «Var. Reduct.» показывает \textbf{истинный выигрыш в точности при строго одинаковом времени работы процессора}.

\begin{lstlisting}
inline MCResult mc_antithetic_variates(const MathFunc& f, double a, double b,
                                       uint64_t n_samples, double exact_val, RandomGenerator& rng) {
    auto t_start = std::chrono::high_resolution_clock::now();
    const double length = b - a;
    const uint64_t n_pairs = n_samples / 2;

    WelfordAccumulator acc;
    for (uint64_t i = 0; i < n_pairs; ++i) {
        double u = rng.uniform01();
        double x1 = a + length * u;
        double x2 = a + length * (1.0 - u);
        double pair_val = 0.5 * length * (f(x1) + f(x2));
        acc.update(pair_val);
    }

    auto t_end = std::chrono::high_resolution_clock::now();
    double ms = std::chrono::duration<double, std::milli>(t_end - t_start).count();

    MCResult res;
    res.method_name = "Antithetic Variates";
    res.n_evals = n_pairs * 2; // Exactly matches n_samples
    res.estimate = acc.get_mean();
    res.variance = 2.0 * acc.get_variance(); // Normalized single-eval variance
    res.std_error = acc.get_std_error();    // std_dev(pair) / sqrt(n_pairs)
    res.ci95_lower = res.estimate - acc.get_ci95_halfwidth();
    res.ci95_upper = res.estimate + acc.get_ci95_halfwidth();
    res.abs_error = std::abs(res.estimate - exact_val);
    res.elapsed_ms = ms;
    return res;
}
\end{lstlisting}

\newpage

\section{Хитрость 4. Искусство подбора контрольных функций (Control Variates)}

Метод контрольных переменных использует известную аналитически функцию $h(x) \approx f(x)$. Однако эффективность метода колоссально зависит от того, \textbf{как именно} построена функция $h(x)$.

\subsection{1. Пример для $\int_0^1 e^x dx$: Хорда против ряда Тейлора}
Рассмотрим две различные стратегии подбора $h(x)$:
\begin{enumerate}
    \item \textbf{Стратегия Тейлора:} Разложение в окрестности $x = 0$: $h_{Taylor}(x) = 1 + x$.
    \begin{itemize}
        \item В точке $x = 0$ ошибка нулевая: $e^0 - 1 = 0$.
        \item В точке $x = 1$ ошибка максимальна: $e^1 - (1 + 1) = e - 2 \approx 0.718$.
        \item Дисперсия остатка: $\sigma^2 \approx 0.0436$. Выигрыш: всего в \textbf{5.5 раз}.
    \end{itemize}
    
    \item \textbf{Наша хитрость --- Двухточечная хорда интерполяции:}
    Соединим прямой линией точные значения функции на концах отрезка $(0, 1)$ и $(1, e)$:
    \begin{equation}
        h_{chord}(x) = 1 + (e - 1)x.
    \end{equation}
    Интеграл от хорды находится тривиально:
    \begin{equation}
        I_h = \int_0^1 \big(1 + (e - 1)x\big) dx = 1 + \frac{e - 1}{2} = \frac{e + 1}{2} \approx 1.8591409.
    \end{equation}
    \begin{itemize}
        \item На обоих концах отрезка $x=0$ и $x=1$ разность $e^x - h(x)$ строго равна нулю!
        \item Максимальное отклонение внутри отрезка нигде не превышает $0.09$ (в 8 раз меньше, чем у Тейлора).
        \item Выборочная дисперсия остатка: $s^2 \approx 0.00400$.
        \item \textbf{Реальный выигрыш: снижение дисперсии более чем в 60 раз!}
    \end{itemize}
\end{enumerate}

\subsection{2. Пример для $\int_0^{\pi/2} \sin(x) dx$: Парабола с нулевой производной}
Для функции $\sin(x)$ на $[0, \pi/2]$:
\begin{itemize}
    \item Если взять линейную хорду $h_1(x) = \frac{2}{\pi}x$, дисперсия снижается в 23 раза.
    \item \textbf{Наша хитрость --- Квадратичная парабола, учитывающая касательную:}  
    Заметим, что в точке $x = \pi/2$ производная функции обращается в ноль: $\sin'(\pi/2) = \cos(\pi/2) = 0$.  
    Построим параболу $h_2(x) = A x + B x^2$, удовлетворяющую условиям:
    $$h_2(0) = 0, \quad h_2(\pi/2) = 1, \quad h_2'(\pi/2) = 0.$$
    Получаем коэффициенты $A = \frac{4}{\pi}$, $B = -\frac{4}{\pi^2}$:
    \begin{equation}
        h_2(x) = \frac{4}{\pi} x - \frac{4}{\pi^2} x^2.
    \end{equation}
    Интеграл вычисляется мгновенно:
    \begin{equation}
        I_h = \int_0^{\pi/2} \left(\frac{4}{\pi} x - \frac{4}{\pi^2} x^2\right) dx = \frac{\pi}{3} \approx 1.04719755.
    \end{equation}
    Разность $\sin(x) - h_2(x)$ настолько мала по всей длине отрезка, что дисперсия падает с $0.2339$ до $0.000941$.
    \item \textbf{Результат: дисперсия снижена в 248.5 раз!}
\end{itemize}

\newpage

\section{Хитрость 5. Краевая сингулярность и машинная арифметика в Варианте 7}

В качестве интеграла из Лабораторной работы №5 для 7 варианта был выбран:
\begin{equation}
    I = \int_0^{\sqrt{3}} \sqrt{3 - x^2} dx = \frac{3\pi}{4} \approx 2.356194490192.
\end{equation}

\subsection{1. Ловушка корня из отрицательного нуля}
В математике для $x \in [0, \sqrt{3}]$ подкоренное выражение $3 - x^2 \ge 0$.  
Однако в численных методах на ЭВМ переменная $x$ генерируется как $x = \sqrt{3} \cdot u$, где $u \in [0, 1)$.

Из-за ошибок округления при представлении $\sqrt{3} \approx 1.7320508075688772$, если датчик случайно выдаст число $u$, чрезвычайно близкое к $1.0$, может возникнуть ситуация:
$$x^2 > 3.0 \implies 3.0 - x^2 = -1.11 \times 10^{-16} < 0.$$
Наивный вызов \texttt{std::sqrt(3.0 - x * x)} немедленно вернет значение \texttt{NaN} (Not a Number), которое заразит весь последующий расчет суммы Велфорда, и вся программа выдаст ошибку!

\subsection{Решение в коде: Защитное отсечение (Clamping)}
В теле лямбда-функции в \texttt{lab6.cpp} реализована проверка знака:
\begin{lstlisting}
auto f\_circ = [](double x) {
    double val = 3.0 - x * x;
    return (val > 0.0) ? std::sqrt(val) : 0.0;
};
\end{lstlisting}
Эта простая строка гарантирует полную надежность вычислений даже при $N = 10^8$ испытаний.

\subsection{2. Сопоставление с поведением формулы Симпсона из Лаб №5}
В управляющей программе \texttt{lab6.cpp} реализован блок прямого сопоставления результатов:
\begin{itemize}
    \item Для функции $\sqrt{3 - x^2}$ производная $f'(x) = -\frac{x}{\sqrt{3 - x^2}}$ стремится к $-\infty$ при $x \to \sqrt{3}$.
    \item Классические квадратуры Ньютона--Котеса (Симпсон, трапеции) предполагают гладкость функции ($f \in C^4$). Из-за сингулярности производной фактический порядок точности формулы Симпсона упал с $O(h^4)$ до $O(h^{1.5})$.
    \item В противоположность этому, методу Монте-Карло гладкость не нужна вовсе. Достаточно лишь конечности интеграла $\int_0^{\sqrt{3}} (3 - x^2)dx = 2\sqrt{3} < \infty$. Метод показал идеальную сходимость со скоростью ровно $O(N^{-0.5})$.
\end{itemize}

\newpage

\section{Сводные результаты работы программы}

Ниже приведены реальные консольные протоколы расчетов программы \texttt{lab6.exe}:

\subsection{1. Задача 1.1: $I = \int_0^1 e^x dx = e - 1$ ($N = 10^6$)}
\begin{table}[h!]
\centering
\small
\begin{tabular}{lrrrrrr}
\toprule
\textbf{Метод} & \textbf{Оценка $\hat{I}$} & \textbf{Дисперсия $s^2$} & \textbf{Станд. ош.} & \textbf{Абс. ош.} & \textbf{Выигрыш} & \textbf{Время (мс)} \\
\midrule
Hit-or-Miss (Геометрический) & 1.717801892 & 1.7186 & $1.31 \times 10^{-3}$ & $4.80 \times 10^{-4}$ & $0.14\times$ & 14.1 \\
Sample Mean (Выборочное ср.) & 1.717974810 & 0.2420 & $4.92 \times 10^{-4}$ & $3.07 \times 10^{-4}$ & $1.00\times$ & 9.1 \\
Antithetic (Антитетический) & 1.718232231 & 0.0078 & $8.85 \times 10^{-5}$ & $4.96 \times 10^{-5}$ & \textbf{30.9$\times$} & 5.9 \\
Control Variates (Контрольн.) & 1.718296254 & 0.0040 & $6.33 \times 10^{-5}$ & $1.44 \times 10^{-5}$ & \textbf{60.5$\times$} & 10.1 \\
Stratified (Расслоенная, K=100)& 1.718271673 & 0.000026 & $5.16 \times 10^{-6}$ & $1.02 \times 10^{-5}$ & \textbf{9087$\times$} & 9.3 \\
\bottomrule
\end{tabular}
\caption{Сравнение методов Монте-Карло для экспоненты при фиксированном $N = 10^6$}
\end{table}

\subsection{2. Сходимость по объему выборки $N$ (логарифмическая сетка)}
\begin{table}[h!]
\centering
\small
\begin{tabular}{rrrrr}
\toprule
\textbf{Число точек $N$} & \textbf{Sample Mean Error} & \textbf{Antithetic Error} & \textbf{Control Var Error} & \textbf{Stratified Error} \\
\midrule
$10^2$ & $5.37 \times 10^{-2}$ & $1.24 \times 10^{-2}$ & $6.39 \times 10^{-3}$ & $1.63 \times 10^{-3}$ \\
$10^3$ & $2.75 \times 10^{-2}$ & $2.69 \times 10^{-4}$ & $2.83 \times 10^{-3}$ & $8.76 \times 10^{-5}$ \\
$10^4$ & $3.33 \times 10^{-3}$ & $5.57 \times 10^{-4}$ & $2.22 \times 10^{-4}$ & $5.05 \times 10^{-5}$ \\
$10^5$ & $1.31 \times 10^{-4}$ & $1.78 \times 10^{-4}$ & $2.96 \times 10^{-4}$ & $3.75 \times 10^{-6}$ \\
$10^6$ & $9.15 \times 10^{-4}$ & $2.26 \times 10^{-5}$ & $1.47 \times 10^{-5}$ & $2.27 \times 10^{-6}$ \\
$10^7$ & $1.29 \times 10^{-4}$ & $1.50 \times 10^{-5}$ & $1.52 \times 10^{-6}$ & $1.77 \times 10^{-6}$ \\
\bottomrule
\end{tabular}
\caption{Эволюция абсолютной погрешности в зависимости от объема выборки $N$}
\end{table}

\subsection{3. Прямое сопоставление с квадратурами Лабораторной работы №5}
\begin{table}[h!]
\centering
\small
\begin{tabular}{lrrcc}
\toprule
\textbf{Метод интегрирования} & \textbf{Вызовы $N$} & \textbf{Погрешность} & \textbf{Порядок точности} & \textbf{Время (мс)} \\
\midrule
Средние прямоугольники (Лаб 5) & 100 & $2.58 \times 10^{-4}$ & $O(h^{1.5})$ & $<0.01$ \\
Формула трапеций (Лаб 5) & 101 & $8.82 \times 10^{-4}$ & $O(h^{1.5})$ & $<0.01$ \\
Формула Симпсона (Лаб 5) & 101 & $3.45 \times 10^{-4}$ & $O(h^{1.5})$ \textit{(деградация с $h^4$)} & $<0.01$ \\
Средние прямоугольники (Лаб 5) & 10\,000 & $2.58 \times 10^{-7}$ & $O(h^{1.5})$ & 0.017 \\
Формула Симпсона (Лаб 5) & 10\,001 & $3.44 \times 10^{-7}$ & $O(h^{1.5})$ \textit{(деградация с $h^4$)} & 0.017 \\
\midrule
Монте-Карло: Sample Mean & 1\,000\,000 & $2.40 \times 10^{-4}$ & $O(N^{-0.5})$ \textit{(строго)} & 7.6 \\
Монте-Карло: Антитетический & 10\,000\,000 & $4.07 \times 10^{-5}$ & $O(N^{-0.5})$ \textit{(строго)} & 46.8 \\
\bottomrule
\end{tabular}
\caption{Прямой батл: Квадратуры Ньютона--Котеса против Монте-Карло для $\int_0^{\sqrt{3}} \sqrt{3-x^2} dx$}
\end{table}

\newpage

\section{Заключение}

В ходе разработки программы были успешно применены следующие ключевые приемы:
\begin{enumerate}
    \item \textbf{Численная устойчивость:} Применение однопроходного алгоритма Велфорда позволило полностью ликвидировать проблему катастрофической потери разрядов при вычислении выборочной дисперсии на выборках объема до $10^7$ точек.
    \item \textbf{Качество псевдослучайных последовательностей:} Использование 64-битного генератора Мерсенна (\texttt{std::mt19937\_64}) исключило цикличность и межвыборочные корреляции.
    \item \textbf{Эффективность понижения дисперсии:} Доказано, что правильный подбор контрольных функций (хорда вместо Тейлора, согласование производных в нуле) способен сократить дисперсию более чем в 240 раз, а расслоенная выборка --- почти в 10\,000 раз.
    \item \textbf{Устойчивость к сингулярностям:} Продемонстрировано фундаментальное преимущество метода Монте-Карло: инвариантность к дифференциальным свойствам подынтегральной функции (отсутствие деградации сходимости на сингулярности $\sqrt{3 - x^2}$).
\end{enumerate}

\end{document}
